\documentclass[10pt,letterpaper]{article}
\usepackage[margin=1in]{geometry}
\usepackage[hidelinks=false]{hyperref}
\usepackage{microtype}
\setlength{\parindent}{0pt}
\setlength{\parskip}{1.25ex}
\pagestyle{plain}
\setcounter{page}{2}
\begin{document}
\vspace*{0.15in}
\begin{center}
{\Large\bfseries Proceedings of International Journal of Secondary\\ Computing and Applications Research, Vol 3, Issue 4}\\[2.4ex]
October 30, 2026
\end{center}
\vspace{1.2ex}
{\large\bfseries Letter from the Editor-in-Chief}\par\vspace{1ex}

It is my pleasure to welcome you to this issue of the International Journal of Secondary Computing and Applications Research (IJSCAR). This is our largest issue yet, with roughly double the number of articles in our previous one, and the breadth of student-led computing research gathered here shows just how quickly this community is growing. The work ranges from machine learning for health and medicine, including genomic variant interpretation, rare-disease modeling, and dental imaging, to computer vision for agriculture, with tomato plant disease detection, and for underwater environments; from predictive maintenance for turbofan engines and anomaly detection in high-energy particle physics to brain-inspired learning algorithms; and from nutrition estimation to audits of demographic sensitivity, with large language models at the center of several studies.

Taken together, these papers show a maturing research instinct: a willingness to engage seriously with real data and real constraints, not just idealized benchmarks. Several authors in this issue compared their methods against strong baselines, examined how sensitive their results are to choices such as decision thresholds, random seeds, and evaluation data, and candidly reported the limits of what their experiments can show rather than reporting a single favorable configuration---exactly the kind of self-scrutiny that distinguishes rigorous research from a class project.

At IJSCAR, our mission remains to provide a platform for students and early-career researchers to share innovative ideas, engage with the broader research community, and contribute to the advancement of knowledge. We are continually inspired by the quality of submissions we receive and by the dedication of young scholars pursuing meaningful and impactful research, and we remain committed to nurturing that curiosity issue after issue.

As Editor-in-Chief, I want to thank our authors for the care they put into this work, their mentors and teachers for the guidance behind the scenes, and our reviewers for the time they gave to strengthening each submission. On behalf of the editorial team, I am also grateful to our readers for their continued support in helping foster a community committed to learning, discovery, and scientific progress.

We hope this issue informs and inspires, and we look forward to sharing the next chapter of student computing research with you soon.

\vspace{1ex}
\begin{tabular}[t]{@{}l@{}}
Sincerely,\\
Maria Hwang\\
Editor-in-Chief
\end{tabular}

\vfill
\hrule\vspace{1.2ex}
Volume 3, Issue 4 \hfill October 30, 2026\par
DOI: 10.67149/yhjs2024.5/d32s2mcm \ $|$ \ \url{https://ijscar.org/pubs/vol3-issue4.pdf}\par
\copyright\ 2026 International Journal of Secondary Computing and Applications Research
\end{document}
